\documentclass[a4paper]{article}

\usepackage[spanish]{babel}
\usepackage{graphicx}
\usepackage{amsmath, amssymb}
\usepackage[margin=2cm]{geometry}
\usepackage{fancyhdr}
\usepackage{enumerate}
\usepackage[shortlabels]{enumitem}
\usepackage{parskip}
\usepackage[most]{tcolorbox}
\usepackage[hidelinks]{hyperref}
\usepackage{float}
\usepackage{booktabs}



% cabecera
\pagestyle{fancy}
\fancyhead[l]{Cosmología Moderna}
\fancyhead[c]{Tareas}
\fancyhead[r]{\today}
\fancyfoot[c]{\thepage}
\renewcommand{\headrulewidth}{0.2pt} % linea horizontal

\begin{document}

\section{Tarea I}

Desde \cite{TroddenCarroll2008}

\begin{enumerate}
	\item Derivar \eqref{eq:17}, \eqref{eq:18} en \cite{TroddenCarroll2008}. Asumir el tensor de Einstein dado en \eqref{eq:1.3.127}, \eqref{eq:1.3.128} de \cite{BaumannTripos}
	
	% Ecuación 17: Ecuación de Friedmann
\begin{equation}
H^2 \equiv \left( \frac{\dot{a}}{a} \right)^2 = \frac{8\pi G}{3} \sum_i \rho_i - \frac{k}{a^2} 
\tag{17}
\label{eq:17}
\end{equation}

% Ecuación 18: Ecuación de Evolución
\begin{equation}
\frac{\ddot{a}}{a} + \frac{1}{2} \left( \frac{\dot{a}}{a} \right)^2 = -4\pi G \sum_i p_i - \frac{k}{2a^2}
\tag{18}
\label{eq:18}
\end{equation}

% Componente temporal del tensor de Einstein (1.3.127)
\begin{equation}
G^0_{\phantom{0}0} = 3 \left[ \left( \frac{\dot{a}}{a} \right)^2 + \frac{k}{a^2} \right]
\tag{1.3.127}
\label{eq:1.3.127}
\end{equation}

% Componentes espaciales del tensor de Einstein (1.3.128)
\begin{equation}
G^i_{\phantom{i}j} = \left[ 2 \frac{\ddot{a}}{a} + \left( \frac{\dot{a}}{a} \right)^2 + \frac{k}{a^2} \right] \delta^i_{\phantom{i}j}
\tag{1.3.128}
\label{eq:1.3.128}
\end{equation}

	\item desde la ecuación \eqref{eq:17} y \eqref{eq:18} derivar \eqref{eq:19}
	
	% Ecuación 19: Ecuación de Aceleración
\begin{equation}
\frac{\ddot{a}}{a} = -\frac{4\pi G}{3} \sum_i (\rho_i + 3p_i)
\tag{19}
\label{eq:19}
\end{equation}
	
	\item desde la ecuación \eqref{eq:17} y \eqref{eq:18} derivar \eqref{eq:25}

% Ecuación 25: Ecuación de Conservación de Energía
\begin{equation}
\dot{\rho} + 3H(\rho + p) = 0
\tag{25}
\label{eq:25}
\end{equation}


	\item desde la ecuación \eqref{eq:17} y \eqref{eq:25} derivar \eqref{eq:18}
\end{enumerate}


\section{Tarea II}

\begin{enumerate}
	\item Mostrar que la ecuación de conservación de la energía para un universo en expansión \eqref{eq:25}, es equivalente a la primera ley de la termodinámica $dE=-PdV$
	\item Derivar \eqref{eq:27} desde la \eqref{eq:25}
	
	% Ecuación 27: Ecuación de Friedmann (Referenciada en la sección 2.3 de Trodden & Carroll)
\begin{equation}
H^2 \equiv \left( \frac{\dot{a}}{a} \right)^2 = \frac{8\pi G}{3} \sum_i \rho_i - \frac{k}{a^2} 
\tag{27}
\label{eq:27}
\end{equation}

	
	\item Derivar \eqref{eq:35}, Asumir $k=\Lambda=0, p=w\rho,a=a_0t^n$ y $\rho(a)$ del punto anterior.

% Ecuación 35: Solución del factor de escala para un universo plano
\begin{equation}
a(t) = a_0 \left( \frac{t}{t_0} \right)^{\frac{2}{3(1+w)}}
\tag{35}
\label{eq:35}
\end{equation}


	\item Mostrar que es el problema del horizonte en universo dominado por materia o radiación
\end{enumerate}



\section{Tarea III}

\begin{enumerate}
	\item Derivar las ecuaciones \eqref{eq:38} hasta \eqref{eq:45}
	
	
	% Ecuación 38: Ecuaciones de Einstein en tiempo conformal (Materia)
\begin{equation}
\begin{aligned}
3(k + h^2) &= 8\pi G \rho a^2 \\
k + h^2 + 2h' &= 0
\end{aligned}
\tag{38}
\label{eq:38}
\end{equation}



% Ecuación 39: Soluciones para h(tau)
\begin{equation}
h(\tau) = \begin{cases} 
\cot(\tau/2) & k = 1 \\ 
2/\tau & k = 0 \\ 
\coth(\tau/2) & k = -1 
\end{cases}
\tag{39}
\label{eq:39}
\end{equation}

% Ecuación 40: Soluciones para a(tau)
\begin{equation}
a(\tau) \propto \begin{cases} 
1 - \cos(\tau) & k = 1 \\ 
\tau^2/2 & k = 0 \\ 
\cosh(\tau) - 1 & k = -1 
\end{cases}
\tag{40}
\label{eq:40}
\end{equation}

% Ecuación 41: Relación entre tiempo cósmico y conformal
\begin{equation}
t(\tau) \propto \begin{cases} 
\tau - \sin(\tau) & k = 1 \\ 
\tau^3/6 & k = 0 \\ 
\sinh(\tau) - \tau & k = -1 
\end{cases}
\tag{41}
\label{eq:41}
\end{equation}

% Ecuación 42: Ecuaciones de Einstein en tiempo conformal (Radiación)
\begin{equation}
\begin{aligned}
3(k + h^2) &= 8\pi G \rho a^2 \\
k + h^2 + 2h' &= -\frac{8\pi G \rho}{3} a^2
\end{aligned}
\tag{42}
\label{eq:42}
\end{equation}

% Ecuación 43: Soluciones para h(tau) (Radiación)
\begin{equation}
h(\tau) = \begin{cases} 
\cot(\tau) & k = 1 \\ 
1/\tau & k = 0 \\ 
\coth(\tau) & k = -1 
\end{cases}
\tag{43}
\label{eq:43}
\end{equation}

% Ecuación 44: Soluciones para a(tau) (Radiación)
\begin{equation}
a(\tau) \propto \begin{cases} 
\sin(\tau) & k = 1 \\ 
\tau & k = 0 \\ 
\sinh(\tau) & k = -1 
\end{cases}
\tag{44}
\label{eq:44}
\end{equation}

% Ecuación 45: Relación t(tau) (Radiación)
\begin{equation}
t(\tau) \propto \begin{cases} 
1 - \cos(\tau) & k = 1 \\ 
\tau^2/2 & k = 0 \\ 
\cosh(\tau) - 1 & k = -1 
\end{cases}
\tag{45}
\label{eq:45}
\end{equation}
	
	\item Derivar \ref{tab:2} por columnas $1$ y $3$
	
	\begin{table}[h]
\centering
\caption{Densidad de número, densidad de energía y presión para especies en equilibrio térmico.}
\label{tab:2}
\begin{tabular}{lccc}
\toprule
 & \textbf{Relativistas} & \textbf{Relativistas} & \textbf{No-relativistas} \\
 & \textbf{Bosones} & \textbf{Fermiones} & \textbf{(Cualquiera)} \\ \midrule
$n_i$ & $\frac{\zeta(3)}{\pi^2} g_i T^3$ & $\left(\frac{3}{4}\right) \frac{\zeta(3)}{\pi^2} g_i T^3$ & $g_i \left(\frac{m_i T}{2\pi}\right)^{3/2} e^{-m_i/T}$ \\ [2ex]
$\rho_i$ & $\frac{\pi^2}{30} g_i T^4$ & $\left(\frac{7}{8}\right) \frac{\pi^2}{30} g_i T^4$ & $m_i n_i$ \\ [2ex]
$p_i$ & $\frac{1}{3} \rho_i$ & $\frac{1}{3} \rho_i$ & $n_i T \ll \rho_i$ \\ \bottomrule
\end{tabular}
\end{table}

\end{enumerate}

\section{Tarea IV}

Desde \cite{BaumannTASI2024}

\begin{enumerate}
	\item Derivar \eqref{eq:A61} y \eqref{eq:A60}
	% Ecuación A.60: Curvatura en densidad uniforme
\begin{equation}
-\zeta \equiv \Psi + H \frac{\delta\rho}{\dot{\bar{\rho}}}
\tag{A.60}
\label{eq:A60}
\end{equation}
% Ecuación A.61: Curvatura comóvil
\begin{equation}
\mathcal{R} = \Psi - H \frac{\delta q}{\bar{\rho} + \bar{p}}
\tag{A.61}
\label{eq:A61}
\end{equation}

	\item Mostrar que la curvatura en hipersuperficies a densidad constante \eqref{eq:A60} y comóviles \eqref{eq:A61} son invariantes de gauge.
	\item Mostrar que las variables de Bardeen \eqref{eq:A53} y \eqref{eq:A54} son invariantes de gauge.
	
	% Ecuación A.53: Primer Potencial de Bardeen
\begin{equation}
\Phi_B \equiv \Phi - \frac{d}{dt} [a^2(\dot{E} - B/a)]
\tag{A.53}
\label{eq:A53}
\end{equation}


% Ecuación A.54: Segundo Potencial de Bardeen
\begin{equation}
\Psi_B \equiv \Psi + a^2 H(\dot{E} - B/a)
\tag{A.54}
\label{eq:A54}
\end{equation}

\end{enumerate}


\section{Tarea V}
Desde \cite{BaumannTripos}
\begin{enumerate}
	\item Derivar \eqref{eq:4.1.24}, \eqref{eq:4.1.25}, \eqref{eq:4.1.26}
	
% Ecuación de Continuidad (4.1.24)
\begin{equation}
\dot{\delta} = -\frac{1}{a} \nabla \cdot \mathbf{v}
\tag{4.1.24}
\label{eq:4.1.24}
\end{equation}

% Ecuación de Euler (4.1.25)
\begin{equation}
\dot{\mathbf{v}} + H\mathbf{v} = - \frac{1}{a\bar{\rho}} \nabla \delta P - \frac{1}{a} \nabla \delta \Phi
\tag{4.1.25}
\label{eq:4.1.25}
\end{equation}

% Ecuación de Poisson (4.1.26)
\begin{equation}
\nabla^2 \delta \Phi = 4\pi G a^2 \bar{\rho} \delta
\tag{4.1.26}
\label{eq:4.1.26}
\end{equation}


	
	\item Derivar \eqref{eq:4.1.30}, \eqref{eq:4.1.34}, \eqref{eq:4.1.38}
	
% Era de Dominación de Materia (MD)
\begin{equation}
\delta_m \propto \begin{cases} t^{-1} \propto a^{-3/2} \\ t^{2/3} \propto a \end{cases}
\tag{4.1.30}
\label{eq:4.1.30}
\end{equation}

% Era de Dominación de Radiación (RD)
\begin{equation}
\delta_m \propto \begin{cases} \text{const.} \\ \ln t \propto \ln a \end{cases}
\tag{4.1.34}
\label{eq:4.1.34}
\end{equation}

% Era de Dominación de Energía Oscura (Lambda)
\begin{equation}
\delta_m \propto \begin{cases} \text{const.} \\ e^{-2Ht} \propto a^{-2} \end{cases}
\tag{4.1.38}
\label{eq:4.1.38}
\end{equation}

	\item Mostrar que \eqref{eq:4.2.120}, \eqref{eq:4.2.115} en el límite no relativista se reducen a las ecuaciones en el punto 1
	
% Ecuación de Continuidad Relativista (Newtonian Gauge)
\begin{equation}
\delta' + \left( 1 + \frac{\bar{P}}{\bar{\rho}} \right) (\nabla \cdot \mathbf{v} - 3\Phi') + 3\mathcal{H} \left( \frac{\delta P}{\delta \rho} - \frac{\bar{P}}{\bar{\rho}} \right) \delta = 0
\tag{4.2.115}
\label{eq:4.2.115}
\end{equation}

% Ecuación de Euler Relativista (Newtonian Gauge)
\begin{equation}
\mathbf{v}' + \mathcal{H}\mathbf{v} - 3\mathcal{H} \frac{\bar{P}'}{\bar{\rho}'} \mathbf{v} = - \frac{\nabla \delta P}{\bar{\rho} + \bar{P}} - \nabla \Psi
\tag{4.2.120}
\label{eq:4.2.120}
\end{equation}
	
\end{enumerate}

\section{Textos}

\begin{itemize}

	\item \cite{dInvernoVickers2022}
	\item \cite{Riotto2002}
	\item \cite{Naruko2018}

\end{itemize}


\bibliographystyle{plainurl}
\bibliography{bibliografia} 
\end{document}
